% 1.5 — Méthodologie
\section{Méthodologie de travail}
\label{sec:methodologie}

\subsection{Le risque propre à un outil de sécurité}
Faire évoluer un outil de sécurité comporte un risque particulier : une modification qui
corrige un cas peut, sans bruit, en autoriser un autre. Un test « vert » ne le révèle pas
si aucun test ne décrivait le cas devenu permissif. Un audit mené au cours du stage l'a montré :
la suite passait avec 1856 tests réussis alors que 18 défauts confirmés restaient
atteignables, parce qu'aucun test ne décrivait ces comportements. La méthode adoptée
répond directement à ce constat.

\subsection{Une partie à la fois}
Le travail est découpé en \emph{parties} courtes et nommées (A, B, C1, SYS1, M2.2b3,
W-SEC\ldots). Chaque partie suit le même cycle (figure~\ref{fig:methode}) :

\begin{enumerate}
  \item \textbf{Plan} : périmètre, cas à couvrir, risques, ce qui ne doit pas bouger.
  \item \textbf{Arbitrages} : chaque question de politique est tranchée \emph{par écrit}
    avant d'être codée, avec un identifiant et une date (« ruling R2 : un envoi qui peut
    remplacer un objet distant est escaladé »). Le code et ses tests citent l'arbitrage.
  \item \textbf{Prototype} : le code, et des tests qui \emph{échouent sur le code d'avant}.
    Un test qui passerait aussi sur l'ancien code ne prouve rien sur la correction.
  \item \textbf{Revue} : relecture adverse et sondes --- des commandes fabriquées pour
    piéger la nouvelle logique, lancées dans un projet jetable.
  \item \textbf{Atterrissage} : un commit dont le message est un bilan : ce qui change, les
    arbitrages appliqués, les tests ajoutés, et le résultat du portail de décision.
\end{enumerate}

\begin{figure}[htbp]
\centering
% Figure 6 — La méthode : une partie à la fois, et les portes qu'elle doit franchir.
\begin{tikzpicture}[
    etape/.style={accent,text width=24mm,minimum height=15mm,font=\sffamily\small},
    porte/.style={boite,draw=encre,text width=30mm,font=\sffamily\scriptsize,minimum height=9mm},
  ]
  \node[etape] (p) at (0,0) {Plan\sous{périmètre, cas, risques}};
  \node[etape,right=6mm of p] (r) {Arbitrages\sous{décisions écrites et datées (R1, M3\ldots)}};
  \node[etape,right=6mm of r] (pr) {Prototype\sous{code + tests qui échouent sur l'ancien code}};
  \node[etape,right=6mm of pr] (rv) {Revue\sous{lecture adverse, sondes}};
  \node[etape,right=6mm of rv,fill=wepwn!35] (l) {Atterrissage\sous{un commit, message = bilan}};
  \foreach \a/\b in {p/r,r/pr,pr/rv,rv/l} \draw[fleche] (\a) -- (\b);
  \draw[fleche,densely dashed] (rv.north) -- ++(0,0.35) -| node[etiquette,pos=0.25,above]{correction} (pr.north);

  % Portes
  \node[porte] (g1) at (2.0,-2.8) {Tests de la partie\sous{rouges avant, verts après}};
  \node[porte,right=4mm of g1] (g2) {Suite complète\sous{Linux \emph{et} Windows}};
  \node[porte,right=4mm of g2] (g3) {Portail de décision\sous{\nbcorpus{} commandes, avant/après}};
  \node[porte,right=4mm of g3] (g4) {Aucune décision plus lâche\sous{sans arbitrage explicite}};
  \node[zone,fit=(g1)(g4),label={[etiquette]below:portes à franchir avant l'atterrissage}] (gz) {};
  \draw[fleche] (gz.north -| l.south) -- (l.south);

  % Boucle de parties
  \node[etiquette,text width=120mm,anchor=north] at ([yshift=-7mm]gz.south) {Parties successives (extrait) :
     A $\cdot$ B $\cdot$ C1 $\cdot$ C2a $\cdot$ SYS1 $\cdot$ S2+M1 $\cdot$ M2.0--M2.1 $\cdot$ T $\cdot$ T5 $\cdot$ B0 $\cdot$ B1 $\cdot$ I1 $\cdot$ W0 $\cdot$
     M2.2a--f $\cdot$ ASK $\cdot$ W-SEC $\cdot$ V17 $\cdot$ M2.3a $\cdot$ M2.3b $\cdot$ V22 $\cdot$ V18 $\cdot$ C3};
\end{tikzpicture}

\caption{La méthode : le cycle d'une partie et les portes à franchir avant l'atterrissage}
\label{fig:methode}
\end{figure}

\subsection{Les portes d'un atterrissage}
\begin{description}[style=nextline,leftmargin=0.6cm,font=\sffamily\bfseries]
  \item[Les tests de la partie] Rouges sur le code d'avant, verts après.
  \item[La suite complète, sous Linux et sous Windows] Les deux systèmes n'ont pas les mêmes
    chemins (\cmd{C:\textbackslash}, \cmd{/c/} sous Git Bash), ni les mêmes appels système, ni le
    même shell : un test vert sous Linux ne dit rien de Windows.
  \item[Le portail de décision] Un corpus de \nbcorpus{} commandes réelles et fabriquées est
    évalué avant et après la modification ; toute ligne dont la décision, la règle, la
    cible ou la capture change est listée et doit être justifiée. Les messages de commit en
    rendent compte (« the 922 gate changes 124 reviewed decisions: 8 fixed misses,
    18 newly captured, 86 newly allowed »).
  \item[Aucune décision plus permissive sans arbitrage] Une ligne qui passe d'ESCALATE à ALLOW
    sans arbitrage explicite bloque l'atterrissage, même si tous les tests passent.
\end{description}

\begin{pourquoi}[Pourquoi un portail de décision en plus des tests ?]
Les tests vérifient ce qu'on a pensé à écrire. Le portail vérifie \emph{tout le reste} :
il ne sait pas quelle décision est juste, mais il voit chaque décision qui change. C'est
lui qui transforme « la suite passe » en « rien n'a bougé sans qu'on le sache ».
\end{pourquoi}

\subsection{Le message de commit comme compte rendu}
Chaque atterrissage laisse un compte rendu lisible directement dans l'historique.
L'extrait suivant est réel (les lignes « [...] » sont coupées) :

\begin{terminal}[git show -s 3c94414]
$ git show -s --format=%B 3c94414
M2.3a: lines the engine does not fully read leave the old text rules

- runners (python, node, npm, make, ./x, source, awk programs that run or
  write ...) go to the runtime net, else a checkpoint, else ALLOW with a
  note that nothing was captured (ruling M1)
- a line the engine cannot read goes to the net, else a checkpoint, else
  ESCALATE (M3) - unread SQL such as psql -c "$(cat q.sql)" was ALLOWED
  (V19); where the old rules flag the line they keep their capture or
  escalation, never an ALLOW
[...]
- engine rules for gzip / bzip2 / xz, sort -o, iconv -o, patch,
  awk -i inplace, 7z: they deleted or overwrote files unseen (M4, V20)
[...]
- tests/test_m23a.py (17); row fixtures of earlier parts: every moved row
  stricter; gate: 6 rows move, reference re-recorded
\end{terminal}
